% from ega2-1-fr.tex
\oldpage[II]{107}

(\hyperref[II.5.3.6-fr]{5.3.6})~; de même, les morphismes $X'_i\to X$ sont
projectifs par (\hyperref[II.5.5.5-fr]{5.5.5, (i) et (ii)}), donc $f$ est
projectif (\hyperref[II.5.5.6-fr]{5.5.6}), et est évidemment surjectif par
définition.

B) Supposons maintenant $X$ irréductible. Comme le morphisme structural
$r:X\to S$ est de type fini, il y a un recouvrement fini $(S_i)$ de $S$ par
des ouverts affines, et pour chaque $i$ un recouvrement fini $(T_{ij})$ de
$r^{-1}(S_i)$ par des ouverts affines, le morphisme $T_{ij}\to S_i$ étant
affine et de type fini, donc quasi-projectif
(\hyperref[II.5.3.4-fr]{5.3.4, (i)})~; comme dans chacune des hypothèses a),
b), l'immersion $S_i\to S$ est quasi-compacte, c'est un morphisme
quasi-projectif (\hyperref[II.5.3.4-fr]{5.3.4, (ii)}), donc la restriction de
$r$ à $T_{ij}$ est un morphisme quasi-projectif
(\hyperref[II.5.3.4-fr]{5.3.4, (iii)}). Désignons les $T_{ij}$ par $U_k$
$(1\leq k\leq n)$. Il existe, pour chaque indice $k$ une immersion ouverte
$\varphi_k:U_k\to P_k$, où $P_k$ est projectif sur $S$
(\hyperref[II.5.3.2-fr]{5.3.2} et \hyperref[II.5.5.2-fr]{5.5.2}). Soit
$U=\bigcap_kU_k$~; comme $X$ est irréductible et les $U_k$ non vides, $U$ est
non vide, et par suite partout dense dans $X$~; les restrictions à $U$ des
$\varphi_k$ définissent un morphisme
\[
\varphi:U\to P=P_1\times_SP_2\times_S\cdots\times_SP_n
\]
de sorte que les diagrammes
\[
\tag{5.6.1.1}
\label{II.5.6.1.1-fr}
\xymatrix{
U\ar[r]^\varphi\ar[d]_{j_k} & P\ar[d]^{p_k}\\
U_k\ar[r]_{\varphi_k} & P_k
}
\]
soient commutatifs, $j_k$ étant l'injection canonique $U\to U_k$ et $p_k$ la
projection canonique $P\to P_k$. Si $j$ est l'injection canonique $U\to X$,
le morphisme $\psi=(j,\varphi)_S:U\to X\times_SP$ est une immersion
(\hyperref[I.5.3.14-fr]{I, 5.3.14}). Dans l'hypothèse a), $X\times_SP$ est
localement noethérien (\hyperref[II.3.4.1-fr]{3.4.1} et
\hyperref[I.6.3.7-fr]{I, 6.3.7 et 6.3.8})~; dans l'hypothèse b),
$X\times_SP$ est un schéma quasi-compact
(\hyperref[I.5.5.1-fr]{I, 5.5.1 et 6.6.4})~; dans les deux cas l'adhérence
$X'$ dans $X\times_SP$ du sous-préschéma $Z$ associé à $\psi$ (et dont
$\psi(U)$ est l'espace sous-jacent) existe, et $\psi$ se factorise en
\[
\tag{5.6.1.2}
\label{II.5.6.1.2-fr}
\psi:U\xrightarrow{\psi'}X'\xrightarrow{h}X\times_SP
\]
où $\psi'$ est une immersion ouverte et $h$ une immersion fermée
(\hyperref[I.9.5.10-fr]{I, 9.5.10}). Soient
$q_1:X\times_SP\to X$ et $q_2:X\times_SP\to P$ les projections canoniques~;
nous poserons
\[
\tag{5.6.1.3}
\label{II.5.6.1.3-fr}
f:X'\xrightarrow{h}X\times_SP\xrightarrow{q_1}X
\]
et
\[
\tag{5.6.1.4}
\label{II.5.6.1.4-fr}
g:X'\xrightarrow{h}X\times_SP\xrightarrow{q_2}P.
\]
Nous allons voir que $X'$ et $f$ répondent à la question.

C) Montrons d'abord que $f$ est projectif et surjectif, et que la restriction
de $f$ à $U'=f^{-1}(U)$ est un isomorphisme de $U'$ sur $U$. Comme les $P_k$
sont projectifs sur $S$, il en est de même de $P$
(\hyperref[II.5.5.5-fr]{5.5.5, (iv)}), donc $X\times_SP$ est projectif sur
$X$ (\hyperref[II.5.5.5-fr]{5.5.5, (iii)}), et il en est de même de $X'$ qui
est un sous-préschéma fermé de $X\times_SP$. D'autre part, on a
$f\circ\psi'=q_1\circ(h\circ\psi')=q_1\circ\psi=j$, donc $f(X')$ contient
l'ensemble ouvert partout dense $U$ de $X$~; mais $f$ est un morphisme fermé
(\hyperref[II.5.5.3-fr]{5.5.3}), donc $f(X')=X$. Notons maintenant que
$q_1^{-1}(U)=U\times_SP$ est induit sur un ouvert de $X\times_SP$, et par
définition le préschéma $U'=h^{-1}(U\times_SP)$ est induit par $X'$ sur
l'ouvert $U'$~; c'est par suite l'adhérence par

